\documentclass[a4paper, 11pt]{article}

% --- Packages ---
\usepackage[utf8]{inputenc}
\usepackage[T1]{fontenc}
\usepackage[french]{babel}
\usepackage[margin=1.5cm, top=2cm, bottom=2cm]{geometry}
\usepackage{xcolor}
\usepackage{fancyhdr}
\usepackage{titlesec}
\usepackage{enumitem}
\usepackage{amsmath}
\usepackage{graphicx}
\usepackage{hyperref}

% --- Définition des couleurs ---
\definecolor{primaryBlue}{RGB}{20, 50, 90}
\definecolor{lightBlue}{RGB}{235, 240, 247}

% --- Styles de document ---
\pagestyle{fancy}
\fancyhf{}
\fancyfoot[L]{\footnotesize \textcolor{gray}{Projet Wakala - Moteur de Recommandation}}
\fancyfoot[R]{\footnotesize \textcolor{gray}{Page \thepage}}
\renewcommand{\headrulewidth}{0pt}
\renewcommand{\footrulewidth}{0.4pt}

% --- Personnalisation des titres ---
\titleformat{\section}{\Large\bfseries\color{primaryBlue}}{}{0em}{}[\titlerule]
\titleformat{\subsection}{\large\bfseries\color{primaryBlue}}{}{0em}{}

\begin{document}

% --- En-tête ---
\begin{center}
    {\Huge \textcolor{primaryBlue}{\textbf{WAKALA}}}\\[0.2cm]
    {\large \bfseries \textcolor{primaryBlue}{Plongée au Cœur du Moteur de Recommandation Hybride}}
\end{center}

\vspace{0.5cm}

\section*{Introduction}
Le moteur de recommandation de la plateforme \textit{Vente-Auto-Platform} est conçu pour offrir aux acheteurs les véhicules les plus pertinents par rapport à leurs besoins et comportements. Contrairement à une simple recherche par filtres, notre système s'appuie sur une approche \textbf{hybride} qui combine la similarité des caractéristiques des véhicules (Content-Based) et les comportements des utilisateurs (Filtrage Collaboratif via Graphe).

\vspace{0.3cm}

\section*{Composants du Moteur de Recommandation}

L'architecture du moteur repose sur quatre piliers fondamentaux :
\begin{itemize}[label=\textbullet]
    \item \textbf{MatchmakerNLP} : Extraction sémantique des intentions utilisateur.
    \item \textbf{Filtrage Basé sur le Contenu (Content-Based)} : Similarité Cosinus mathématique.
    \item \textbf{Filtrage Collaboratif (Neo4j)} : Analyse comportementale en graphe.
    \item \textbf{Moteur Hybride (HybridEngine)} : Agrégation intelligente avec A/B testing.
\end{itemize}

\vspace{0.2cm}

\subsection*{1. Matchmaker NLP : Comprendre l'Utilisateur}
Avant de chercher un véhicule, le système doit comprendre des requêtes en langage naturel telles que \textit{"Je cherche une voiture familiale spacieuse pour moins de 200 000 DH"}.
Le composant \texttt{MatchmakerNLP} convertit ce texte libre en un dictionnaire de caractéristiques structurées.

\begin{itemize}[label=\textbullet]
    \item \textbf{Extraction Budgétaire :} Utilisation d'expressions régulières (Regex) pour capturer les montants numériques associés aux mots "budget", "moins de" ou "max".
    \item \textbf{Classification de Style de vie (Lifestyle) :} Détection de mots-clés :
    \begin{itemize}
        \item \textit{Famille} ("enfants", "minivan", "spacieux") $\rightarrow$ Filtres : SUV, Break.
        \item \textit{Urbain} ("ville", "citadine") $\rightarrow$ Filtres : Citadine, Compacte.
        \item \textit{Écologique} ("hybride", "électrique").
        \item \textit{Sportif} ("rapide", "puissant").
    \end{itemize}
\end{itemize}

\vspace{0.2cm}

\subsection*{2. Filtrage Basé sur le Contenu (Content-Based)}
Le composant \texttt{content\_based.py} évalue mathématiquement à quel point un véhicule correspond aux critères de l'utilisateur.

\textbf{Vectorisation et Normalisation}\\
Chaque véhicule est transformé en un vecteur de 7 dimensions (les \textit{Features}) :
\begin{equation}
    V_i = [Prix, Ann\acute{e}e, Kilom\acute{e}trage, Type, Bo\hat{i}te, Carburant, Puissance]
\end{equation}

Les données ayant des échelles très différentes, le système applique des algorithmes de normalisation (\texttt{StandardScaler} et \texttt{MinMaxScaler} de \textit{scikit-learn}) pour ramener toutes les valeurs sur des échelles comparables.

\vspace{0.2cm}

\textbf{Similarité Cosinus}\\
Le système construit un vecteur requête $Q$ basé sur les critères de l'utilisateur, et calcule la \textbf{Similarité Cosinus} entre $Q$ et chaque véhicule $V_i$ :
\begin{equation}
    \text{Similarité}(Q, V_i) = \frac{Q \cdot V_i}{\|Q\| \times \|V_i\|}
\end{equation}
Le résultat est normalisé entre 0 et 1 pour générer le \textbf{Content Score}.

\vspace{0.2cm}

\subsection*{3. Filtrage Collaboratif (Neo4j)}
L'approche basée sur le contenu ne suffit pas toujours. Si un utilisateur cherche un SUV, il est pertinent de lui recommander les SUV que les autres utilisateurs au profil similaire ont consultés ou mis en favoris.

La base de données \textbf{Neo4j} stocke les entités (Users, Vehicles) sous forme de Nœuds, et les interactions sous forme de Relations (\texttt{VIEWED}, \texttt{FAVORITED}).
En interrogeant le graphe, le système génère un \textbf{Collaborative Score} mesurant la popularité comportementale d'un véhicule par rapport à la "bulle" de l'utilisateur.

\vspace{0.2cm}

\subsection*{4. Le Moteur Hybride et l'A/B Testing}
Le \texttt{HybridEngine} fusionne les résultats des deux méthodes précédentes via une pondération :
\begin{equation}
    \text{Score Final} = \alpha \times \text{ContentScore} + (1 - \alpha) \times \text{CollaborativeScore}
\end{equation}

Pour optimiser l'algorithme, le système effectue de l'A/B Testing en assignant aléatoirement les utilisateurs connectés dans deux seaux :
\begin{itemize}[label=\textbullet]
    \item \textbf{Variante A ($\alpha = 0.8$)} : Privilégie massivement le moteur \textit{Content-Based} (caractéristiques techniques du véhicule).
    \item \textbf{Variante B ($\alpha = 0.2$)} : Privilégie massivement le moteur \textit{Collaboratif} (comportement social et tendances).
\end{itemize}

\vspace{0.3cm}

\section*{Conclusion}
En combinant le NLP, la Similarité Cosinus vectorielle et la puissance de la théorie des Graphes, le Moteur de Recommandation de la \textit{Vente-Auto-Platform} assure une personnalisation de pointe, évolutive et auto-apprenante grâce à son infrastructure d'A/B Testing native.

\end{document}
